Put \(J=D\cap\Delta_s\) and \(I=D\setminus\Delta_s\).  The displayed
partial-selection hypothesis gives
\[
J\ne\varnothing,\qquad I\ne\varnothing,\qquad D\cap\Delta_s\ne\varnothing.
\tag{1}
\]
Consequently the antecedent \(D\cap\Delta_s=\varnothing\) in
\texttt{ass:district-kernel-invariance} is false for the pair \((s,D)\).
After district closure is omitted, that assumption therefore imposes no
equation between \(q_D^s\) and \(q_D^\star\); in particular, it implies
neither equality of their full district response laws nor equality of any
proper coordinate marginal.

Let \(\ell\in\{0,1\}^{\Omega_D}\) be a nonconstant incidence row.  By
nonconstancy, there are \(\omega_0,\omega_1\in\Omega_D\) such that
\[
\ell(\omega_0)=0,\qquad \ell(\omega_1)=1.
\tag{2}
\]
For \(i\in\{0,1\}\), construct a source-labelled family
\(\mathbf M^{(i)}\) by assigning the focal-district response laws
\[
q_D^{s,(i)}=\delta_{\omega_0},
\qquad
q_D^{\star,(i)}=\delta_{\omega_i}.
\tag{3}
\]
For every domain \(t\in\mathcal S\setminus\{s,\star\}\), assign
\[
q_D^{t,(i)}=
\begin{cases}
\delta_{\omega_i},&D\cap\Delta_t=\varnothing,\\
\delta_{\omega_0},&D\cap\Delta_t\ne\varnothing.
\end{cases}
\tag{4}
\]
For every district \(D'\ne D\), choose a deterministic response profile
\(\bar\omega_{D'}\) and use the point mass
\(\delta_{\bar\omega_{D'}}\) in every domain and in both families.

We verify that these response laws are realized by finite SCMs on the
declared alphabets.  By \texttt{def:district-response-profile}, each
\(\omega_i\) is a tuple of deterministic structural response maps taking
values in the declared state spaces.  At every parent configuration not
visited by the target or supplied regimes, extend each such map by a fixed
element \(x_V^\circ\in\mathcal X_V\).  Such an element exists because each
declared state space is nonempty.  Use the completed maps as the structural
tables, allowing all graph-permitted directed-parent arguments, and let the
associated exogenous variables have one-point laws selecting the required
tables.  Because \(G\) is acyclic, recursive evaluation in a topological
order is well defined; every table has codomain \(\mathcal X_V\), so the
result is a finite original-alphabet SCM satisfying
\texttt{ass:recursive-compatibility}.  All exogenous response blocks are
degenerate, and hence their joint law is the product of their marginal
laws; thus \texttt{ass:district-factorization} holds.

For every \(t\) with \(D\cap\Delta_t=\varnothing\), equation \((4)\)
sets the law on \(D\) equal to the target law.  In all nonfocal districts,
the same point mass is used in every domain.  These facts verify every
instance of \texttt{ass:district-kernel-invariance} whose antecedent is
true.  Its antecedent is false for \((s,D)\) by \((1)\).  Hence both
\(\mathbf M^{(0)}\) and \(\mathbf M^{(1)}\) satisfy all retained frozen
assumptions.

The SCM in the partially selected source is identical in the two families:
\[
M^{s,(0)}=M^{s,(1)}.
\tag{5}
\]
Therefore every supplied observational, interventional, or declared
partial-copy law evaluated in that source is identical in the two
families.  In particular, in the source-row-only information configuration
of the claim, the supplied row value is
\[
\sum_{\omega\in\Omega_D}\ell(\omega)q_D^{s,(0)}(\omega)
=
\sum_{\omega\in\Omega_D}\ell(\omega)q_D^{s,(1)}(\omega)
=\ell(\omega_0)=0.
\tag{6}
\]
By contrast, the target values are
\[
\sum_{\omega\in\Omega_D}\ell(\omega)q_D^{\star,(0)}(\omega)=0,
\qquad
\sum_{\omega\in\Omega_D}\ell(\omega)q_D^{\star,(1)}(\omega)=1,
\tag{7}
\]
where \((7)\) follows from \((2)\)--\((3)\).  Equations
\((3)\)--\((7)\) give the required original-alphabet paired-model
certificate.  Notice that a supplied law from an additional wholly
unselected domain could convey target-district information by the retained
invariance assumption; such a law is not part of the source-row-only
information configuration asserted here.

Conversely, if \(\ell\equiv c\) is constant, then normalization of every
\(q\in\Delta(\Omega_D)\) gives
\[
\sum_{\omega\in\Omega_D}\ell(\omega)q(\omega)
=c\sum_{\omega\in\Omega_D}q(\omega)=c.
\tag{8}
\]
Thus constant rows transport universally, while \((3)\)--\((7)\) refute
universal transport of every nonconstant binary row.  This proves the first
characterization.

For the second characterization, let
\[
\pi_I:\Omega_D\longrightarrow\pi_I(\Omega_D)
\]
be coordinate restriction to the response maps of the vertices in \(I\),
and impose the additional marginal-invariance premise
\[
(\pi_I)_\#q_D^s=(\pi_I)_\#q_D^\star.
\tag{9}
\]
Suppose first that \(\ell\) is constant on every fiber of \(\pi_I\).
Then the function
\[
\widetilde\ell(u)=\ell(\omega)
\quad\text{for any }\omega\in\Omega_D\text{ with }\pi_I(\omega)=u
\]
is well defined and satisfies \(\ell=\widetilde\ell\circ\pi_I\).  Grouping
the finite sums by fibers and using \((9)\) yields
\[
\begin{aligned}
\sum_{\omega\in\Omega_D}\ell(\omega)q_D^s(\omega)
&=\sum_{u\in\pi_I(\Omega_D)}
  \widetilde\ell(u)(\pi_I)_\#q_D^s(u)\\
&=\sum_{u\in\pi_I(\Omega_D)}
  \widetilde\ell(u)(\pi_I)_\#q_D^\star(u)\\
&=\sum_{\omega\in\Omega_D}\ell(\omega)q_D^\star(\omega).
\end{aligned}
\tag{10}
\]
Thus fiber constancy is sufficient.

For necessity, suppose that \(\ell\) is not constant on the fibers.  Then
there are \(\omega_0,\omega_1\in\Omega_D\) such that
\[
\pi_I(\omega_0)=\pi_I(\omega_1),
\qquad
\ell(\omega_0)\ne\ell(\omega_1).
\tag{11}
\]
Use these two profiles in the construction \((3)\)--\((4)\).  Extend their
unvisited structural tables so that their coordinates on \(I\) remain
identical.  The resulting source and target laws obey the added premise
because
\[
(\pi_I)_\#\delta_{\omega_0}
=\delta_{\pi_I(\omega_0)}
=\delta_{\pi_I(\omega_1)}
=(\pi_I)_\#\delta_{\omega_1}.
\tag{12}
\]
Nevertheless their row values are \(\ell(\omega_0)\) and
\(\ell(\omega_1)\), which differ by \((11)\).  Hence a row transports for
every pair satisfying \((9)\) only if it is constant on every projection
fiber.  Together with \((10)\), this proves the asserted if-and-only-if.

The condition is decidable because \(\Omega_D\) is finite.  Enumerate all
pairs \((\omega,\omega')\in\Omega_D^2\), and reject precisely when
\[
\pi_I(\omega)=\pi_I(\omega')
\quad\text{and}\quad
\ell(\omega)\ne\ell(\omega').
\tag{13}
\]
If no pair is returned, exhaustive enumeration proves fiber constancy.  If
a pair is returned, the two point masses at that pair, as in
\((11)\)--\((12)\), are the leakage certificate prescribed by
\texttt{def:partial-district-leakage-handle}.

For the smallest exact stress check, take the binary ADMG
\(X\leftrightarrow Y\), let \(D=\{X,Y\}\), and let
\(\Delta_s=\{X\}\).  The deterministic profiles
\(\omega_0=(X=0,Y=0)\) and \(\omega_1=(X=1,Y=0)\) have the same
\(I=\{Y\}\) coordinate.  For
\(\ell(\omega)=\mathbf 1\{X(\omega)=1\}\), identical source point masses
at \(\omega_0\) and target point masses at \(\omega_0\) and \(\omega_1\)
give target values \(0\) and \(1\), respectively, while preserving the
\(I\)-marginal.  This is the advertised two-point partial-district leakage
witness.  No conditioning, division, matrix inversion, limiting argument,
rank condition, positivity condition, or overlap condition is used in the
proof; the only boundary property used is that point masses belong to the
finite probability simplex.